\section{The uniform logarithmic gain and its actual local-lattice interface}
\label{sec:global-local-gain-v046}

\status{The logarithmic counting theorem in this chapter is a theorem
about every original admitted family, with the original constants and
the error \(b^{x^{-\epsilon/2}}\). The subsequent local-lattice
installation retains the original comparison and growth hypotheses.
Neither a replacement family nor an independent collection of local
lattices is introduced. The whole global local-gain obligation is
closed only if those actual-family hypotheses are supplied.}

The source revision throughout is
\texttt{7cbf7a0c32fca1badf05931879a37c3f8a2b03b4}.
Source declarations specify the objects being reconstructed; the
arguments below supply the mathematical proofs. Previously proved
local projection theorems TC6, EA9, GI8, GII8, BC8 and GJ4 are reused.
In particular the counting argument below does not posit a new local
projection theorem.

\subsection{The original quantities}

Fix \(0<\epsilon<1/52\), an original member \(D\), its size datum \(L\),
and its instance \(X\). Write
\[
 \begin{gathered}
 b=\exp(x^\epsilon),\quad a=xb,\quad
 X'=\exp(x^{\epsilon/3}),\quad b'=a/X',\\
 \delta=x^{-\epsilon/8},\quad
 w=b'/(xn),\quad X''=a/w=xnX',\quad
 p_0=k^2x^{26},\\
 N_m=k+m(t+1),\qquad S_c=4k\Delta,\qquad c=4S_c,
 \qquad \eta=L.\mathrm{etaSz}.
 \end{gathered}                                                    \tag{GLG.1}
\]
Here \(m\ge2\), \(k\ge m+3\), \(n=|HH|\ge1\), \(f=2n\),
and \(t\) is the original positive odd integer. A prime always means
a rational prime. No bound on the functions \(b_s\) from above will
be imposed. The original window is
\[
 a\le a_s\le a+\frac{a}{kx^{13}},\qquad b_s\ge b'
                         \quad(s\in TT).                         \tag{GLG.2}
\]
The \(a_s,b_s\) are integers, although all estimates below are real.

The constants \(K\) are the original nonnegative record
\(\mathrm{Consts}\). To make the formulas readable set
\[
 \begin{gathered}
 K_D=K.\mathrm{DII},\quad K_I=K.\mathrm{DI},\quad
 K_T=K.\mathrm{tII},\quad K_P=K.\mathrm{Pr1},\\
 K_1=K.\mathrm{kappaThresholdOne},\quad
 K_2=K.\mathrm{kappaThresholdTwo}.
 \end{gathered}
\]
All other fields, notably \(K.\mathrm{dl1}\), \(K.\mathrm{dl2}\),
\(K.\mathrm{La}\) and \(K.\mathrm{kII}\), retain their original
meaning. Put
\[
 \begin{split}
 r&=\left\lceil8(4+8\epsilon)/\epsilon\right\rceil+1,\\
 L_1(p)&=\left\lceil\log_p((kx^{26})^4)\right\rceil,\qquad
 \Lambda=K.\mathrm{La}\,N_1^2\log a,\qquad Q=Q_p .
 \end{split}                                                       \tag{GLG.3}
\]
The \(N_1,G_p,L_2\) below are exactly the original orbit-window count,
good-window length and twice the supremum of \(L_1(p)\) over inert
primes. All ceilings of nonnegative real numbers are natural ceilings.

For \(1\le j\le m\), put
\(\kappa_j=[k+(2j-1)(t+1)]/2\).
The original filling function is \(F_j(z)=X.\mathrm{fill}(b,j,z)\):
at a top \((0,s)\) it is \(b_s+o_j(0,s)\), where
\[
 b'/x\le o_j(0,s)\le b'/x+\eta.
\]
At a thin place its offset lies in \([w,w+\eta]\), and at a thick
place in \([b,b+\eta+1]\). Set
\(F_j^G=F_j-\kappa_j\), \(\alpha_s=a_s-2c\), and
\(U_j(s,e)=F_j(\operatorname{partner}(p,s,e))-\kappa_j+c\).
Thus \(F_j,F_j^G,U_j\) are the source's \(\mathrm{fill}\),
\(\mathrm{fillG}\), and \(\mathrm{capU}\).

The additive orbit step is \(h_p\) at inert primes and \(2h_p\)
at signal primes. Let \(R_p\) be the chosen representatives, let
\(T_v\) be the tops with representative \(v\), and let \(N_v=|T_v|\).
The source's \(n_p\) is \(\operatorname{ord}(h_p)\) at inert primes,
and \(2\operatorname{ord}(2h_p)\) at signal primes. The subset \(V\)
of a cyclic orbit consists of its even positions, measured from the
same chosen representative. A \emph{candidate} prime is inert,
\(p<X'\), of order at least \(x^{\epsilon/4}\), and satisfies the
original exclusion-band condition for \(p\ge p_0\). In particular
its order \(n'\) is even. These are the original \(\mathrm{Cand}\)
and \(\mathrm{IsV}\), with the orbit dictionary proved in GJ3.

Write \(u_+=\max(0,u)\) and \(\gamma_p=p/(p^2-1)\). The relevant
local quantities are
\[
 \begin{split}
 A_j(s)&=\sum_{\substack{e\ge1\\ e\ {\rm odd}}}
       (a_s/p^e-F_j(\operatorname{partner}(p,s,e)))_+,\\
 A_j^G(s)&=\sum_{\substack{e\ge1\\ e\ {\rm odd}}}
       (a_s/p^e-F_j^G(\operatorname{partner}(p,s,e)))_+,\\
 M_j(s)&=\sum_{\substack{e\ge1\\ e\ {\rm odd}}}
       \min(a_s/p^e,F_j(\operatorname{partner}(p,s,e))),\\
 d_v&=K.\mathrm{dl1}\,n_pN_m^2\log_p(2a)
          +K.\mathrm{dl2}\,N_v\log_p(2a)^2,\\
 g_v^0&=\sum_{s\in T_v}
       \sum_{\substack{e\ge1\ {\rm odd}\\\neg\mathrm{Excl}(p,s,e)}}
          (\alpha_s/p^e-U_j(s,e))_+-d_v,\\
 g_v^U&=\gamma_p\sum_{s\in T_v}\alpha_s
       -\sum_{s\in T_v\setminus V}\sum_{e\ {\rm odd}}
          \min(\alpha_s/p^e,U_j(s,e))-d_v .
 \end{split}                                                       \tag{GLG.4}
\]
They are \(\mathrm{askP},\mathrm{askG},\mathrm{Mn},\mathrm{dl},
\mathrm{g0},\mathrm{gU}\), respectively. Positive-part sums have
finite support; the minimum sums converge by comparison with
\(\sum a_s/p^e\).

\subsection{The finite counting argument}

% BEGIN GLG1
\paragraph{GLG1. numerical bounds on the original filling data.}
\label{claim:GLG1}
Under the original \(\mathrm{ParameterChoice}'\),
\[
 \begin{gathered}
 k\ge5,\quad N_m\le S_c,\quad S_c\ge100,\quad
 500\le w,\quad kS_c\le w\le b'/x\le b',\quad 2b'\le a,\\
 0\le\kappa_j\le N_m\le S_c,\qquad
 F_j\ge1,\quad U_j\ge0,\quad
 0\le\alpha_s\le6a/5. 
 \end{gathered}                                                     \tag{GLG.5}
\]
Also \(F_j^G\ge1\) whenever \(b_s\ge0\) on \(TT\) and
the original size condition \(N_m+1\le w\) holds.

\emph{Proof.}
The rank conditions give \(k\ge5\) and
\(\Delta\ge N_m m^{10}\ge N_m\ge k\).
Since \(S_c=4k\Delta\), this gives \(N_m\le S_c\) and
\(S_c\ge100\). The parameter inequality \(kS_c\le w\)
then gives \(w\ge500\). The formula for \(w\), \(n\ge1\)
and \(x>1\) give its two upper bounds. Since \(X'>2\),
\(2b'\le a\). The formula for \(\kappa_j\), \(1\le j\le m\)
and \(N_m=k+m(t+1)\) give \(0\le\kappa_j\le N_m\).
Every filling is at least \(w\): a top is at least \(b'/x\),
a thin place is at least \(w\), and a thick place is at least
\(b>w\), because \(w=b/(nX')\). This proves the filling assertions.
Furthermore \(c=4S_c\le w\le b'\), so \(2c\le a\).
Finally \(kx^{13}\ge5\), and (GLG.2) gives
\(a_s\le a+a/5=6a/5\). Subtracting \(2c\) proves the last
claim. All these comparisons hold for the actual parameters;
no asymptotic replacement of \(a_s\) or \(b_s\) has occurred.
\hfill\(\square\)
% END GLG1

% BEGIN GLG2
\paragraph{GLG2. finite shifts of the odd geometric sum.}
\label{claim:GLG2}
Let \(p\ge2\), \(c\ge0\), \(2c\le y\le A\), \(A\ge1\).
Suppose \(W_e\ge1\), \(0\le U_e\le W_e+c\). For any subset
\(E\) of the positive odd integers,
\[
 \sum_{e\in E}(y/p^e-W_e)_+
 \le\sum_{e\in E}((y-2c)/p^e-U_e)_+
                                  +3c\log_p A.                    \tag{GLG.6}
\]
If \(0\le\kappa\) and \(W_e-\kappa\ge1\), then
\[
 0\le\sum_{e\ {\rm odd}}(y/p^e-W_e+\kappa)_+
             -\sum_{e\ {\rm odd}}(y/p^e-W_e)_+
                      \le\kappa\log_p A                           \tag{GLG.7}
\]
for \(0\le y\le A\). Moreover
\[
              A_j(s)+M_j(s)=\gamma_p a_s .                         \tag{GLG.8}
\]

\emph{Proof.}
In (GLG.6), a positive term on the left requires \(p^e<y\le A\).
For each such exponent the difference of the positive parts is at
most \(2c/p^e+c\le3c\). There are at most \(\log_p A\)
positive integer exponents with \(p^e<A\). All other left terms
vanish, while their right terms are nonnegative. This proves
(GLG.6), including \(A<p\), when both contributing sets on the left
are empty. For (GLG.7), the pointwise difference is between zero
and \(\kappa\); it can be nonzero only if
\(y/p^e>W_e-\kappa\ge1\). The same exponent count applies.
For (GLG.8) use
\((u-v)_++\min(u,v)=u\) for \(u,v\ge0\), and sum
\(\sum_{e\ge1,\ e\ {\rm odd}}p^{-e}=p/(p^2-1)\).
The geometric bound also justifies the interchange with every
finite sum over tops.
\hfill\(\square\)
% END GLG2

% BEGIN GLG3
\paragraph{GLG3. the actual cyclic orbit and the nearest even predecessor.}
\label{claim:GLG3}
For a candidate prime and a representative \(v\), the map
\[
 \phi:\mathbb Z/n'\mathbb Z\longrightarrow HH,\qquad
       i\longmapsto v+i h_p,\qquad n'=\operatorname{ord}(h_p),
                                                                    \tag{GLG.9}
\]
is injective, identifies the original tops and even positions,
and sends addition of \(e\) to the original partner
\((0,\phi(i)+e h_p)\). For an odd top \(y\), let \(d(y)\) be
the least positive backward distance to an even top, or zero
if no such top exists. Put \(x(y)=y-d(y)\) when \(d(y)>0\).
For nonnegative \(a_x,W_y\),
\[
 \sum_{\substack{y\ {\rm odd\ top}\\d(y)>0}}
          \min(W_y,a_{x(y)}/p^{d(y)})
 \le\sum_{x\ {\rm even\ top}}\sum_{e\ {\rm odd}}
          \min(a_x/p^e,W_{x+e}).                                  \tag{GLG.10}
\]
Let
\[
 S_{\rm kink}=\sum_{\substack{y\ {\rm odd\ top}\\d(y)>0}}
                      \min(b_y,a_{x(y)}/p^{d(y)}).
 \]
Then, on the actual orbit,
\[
 \sum_{s\in T_v\setminus V}(F_j^G(0,s)-b_s)+S_{\rm kink}
       \le\sum_{s\in T_v\cap V}M_j(s)+N_v(b'/x+\eta).              \tag{GLG.11}
\]

\emph{Proof.}
If \(\phi(i)=\phi(i')\), the order of \(h_p\) divides \(i-i'\),
so the residues coincide. The representative function is constant
on an orbit and fixes its representative. Its fiber is therefore
exactly the image of (GLG.9). This proves the dictionary, including
the cardinal \(N_v\), without replacing the selected representative.
The evenness of \(n'\) makes parity well defined modulo \(n'\);
thus \(d(y)\) is odd. The pair \((x(y),d(y))\) determines
\(y=x(y)+d(y)\), so the assignment in (GLG.10) is injective.
Its terms occur among the nonnegative terms on the right.

For (GLG.11), use \(W_y=F_j(0,y)\) in (GLG.10).
At a top \(F_j(0,y)=b_y+o_j(0,y)\ge b_y\), while
\[
 F_j^G(0,y)-b_y=o_j(0,y)-\kappa_j\le b'/x+\eta.
 \]
Consequently each left summand of (GLG.11) is bounded by
the corresponding summand of (GLG.10) plus \(b'/x+\eta\).
For a top with no predecessor only this last bound is used.
There are at most \(N_v\) such tops. The partner dictionary
identifies the right side of (GLG.10) with the even-top
minimum sums \(M_j\), proving the assertion.
\hfill\(\square\)
% END GLG3

% BEGIN GLG4
\paragraph{GLG4. an explicit numerical budget for one good orbit.}
\label{claim:GLG4}
Set \(l=\log x\), \(l_p=\log_p a\), \(N=N_v\), and
\[
 \begin{gathered}
 G_0=C_{31}l/\delta,\qquad U=w+\eta+c+N_m,\qquad c_0=2b'+2c,\\
 E_1=\Lambda a/(pk)+ap^{-L_1(p)},\qquad
 E_2=a/p^2+\Lambda a/(pk),\\
 W_D=K_D(N_1^3+r^5+k^2)l_p,\qquad
 E_I=K_Ix^{26}(N_1^5+k^3)l_p,\\
 D_I(n')=\lceil40n'\delta E_I\rceil,\qquad
 c_6(n')=8(D_I(n')+1)+40K_2n'\delta k l_p
                       +160n'\delta c+2n'c .
 \end{gathered}                                                     \tag{GLG.12}
\]
Assume the original \(\mathrm{adm1}(C_{31})\), \(C_{31}>0\), and
the following numerical conditions:
\[
 \begin{gathered}
 L_2<\lceil1/\delta\rceil,\quad 3G_p\le x^{\epsilon/4},
 \quad l\ge1,\quad G_0U\le b',\quad c_0\le3b',\\
 kQ_p\le b'\ (p<p_0),\qquad W_D\le b'\ (\text{every prime }p),\\
 C\ge3,\quad C\ge K_Tr,\quad
 C\ge13+4K_T+160C_{31},\quad C\ge K_PC_{31},\\
 K_1N_1\le Cl ,
 \end{gathered}                                                     \tag{GLG.13}
\]
and, for \(p_0\le p<X'\) and every \(n'\le n\),
\[
 \begin{split}
 B(n',p):={}&K_PG_0(\eta+c+N_m)
       +K_2(Q_pc_6(n')+kQ_p+kl_p+c)\\
       &+2G_0c+40C_{31}lE_I+3\le b_o .
 \end{split}                                                       \tag{GLG.14}
\]
These are precisely the original \(\mathrm{NumStar}\) inequalities.
For every good candidate orbit,
\[
              N\ge1,\qquad n'\delta\le NC_{31}l,\qquad n'\le NG_0 .
                                                                    \tag{GLG.15}
\]
Moreover
\[
 \sum_{s\in T_v}A_j(s)+\sum_{s\in T_v\cap V}M_j(s)
                  -3Nc\log_p(2a)\le g_v^U+d_v .                   \tag{GLG.16}
\]

\emph{Proof.}
The order bound and the first two conditions in (GLG.13) activate
the original good-orbit count in \(\mathrm{adm1}\):
\(n'\delta/(C_{31}l)\le N\).
Its positive numerator gives \(N>0\), hence \(N\ge1\);
the other two assertions follow by multiplication.

For an odd-position top, apply (GLG.6) with
\(y=a_s\), \(A=2a\), and \(U_e=U_j(s,e)\). The hypotheses
hold by GLG1, and the geometric identity for \(a_s-2c\)
turns the right side into
\[
 \gamma_p\alpha_s-\sum_{e\ {\rm odd}}
                         \min(\alpha_s/p^e,U_j(s,e)).
 \]
For an even-position top, (GLG.8) gives
\(A_j(s)+M_j(s)=\gamma_p(\alpha_s+2c)\).
As \(p<X'<a\), \(\log_p(2a)\ge1\); also \(\gamma_p\le1\).
Thus \(2c\gamma_p\le3c\log_p(2a)\).
Sum the two kinds of tops and use the definition of \(g_v^U\).
This proves (GLG.16) with its original loss, independently of
the size of any \(b_s\) above its lower bound.
\hfill\(\square\)
% END GLG4

% BEGIN GLG5
\paragraph{GLG5. the good-orbit inequality below \(p_0\).}
\label{claim:GLG5}
Under GLG4, let a good candidate prime satisfy \(p<p_0\).
The original quantities on this orbit are
\[
 \begin{split}
 \Pr_v&=r\{N(2b'+c)+n'U\},\\
 t_v&=K_TrN(E_1+c_0+kQ)+2rn'c,\\
 D_p&=\left\lceil40n'\delta W_D+2rn'c
                                      +40rn'\delta c_0\right\rceil .
 \end{split}                                                       \tag{GLG.17}
\]
For every original lift \(j\),
\[
 \sum_{s\in T_v}\{A_j(s)-\operatorname{lo}(C,b_o,p)\}
       \le g_v^U+d_v-\Pr_v-\lceil t_v\rceil-D_p-1.                 \tag{GLG.18}
\]
This is the full original \(\mathrm{StarAt}\) in this prime range.

\emph{Proof.}
All arguments of the ceilings are nonnegative. From (GLG.15)
and \(G_0U\le b'\), both \(n'U\) and \(n'c\) are at most
\(Nb'\). The deliberately coarse bound
\(\Pr_v\le6rNb'\) follows. The part of \(t_v\) other than
\(K_TrNE_1\) is at most \((4K_T+2)rNb'\).
The argument of the final ceiling is at most
\[
       40C_{31}Nl b'+2rNb'+120rC_{31}Nl b'.
 \]
The two ceilings and the final \(1\) cost at most \(3\).
Since \(N,r,l,b'\ge1\), the sum of these non-\(E_1\) terms
is at most
\[
              (13+4K_T+160C_{31})Nr b'l .
 \]
The remaining \(E_1\) term is at most \(NClE_1\), using
\(K_Tr\le C\) and \(l\ge1\). Consequently
\[
       \Pr_v+\lceil t_v\rceil+D_p+1
                        \le NCr b'l+NClE_1.                     \tag{GLG.19}
\]
In this prime range the original definition is
\[
 \operatorname{lo}(C,b_o,p)
             =C(c+\eta)\log_p(2a)+Crb'l+ClE_1.
 \]
Its first term absorbs \(3c\log_p(2a)\) by \(C\ge3\).
Subtract (GLG.19) from (GLG.16) and discard the nonnegative
even-top sum \(M_j\). This is (GLG.18).
\hfill\(\square\)
% END GLG5

% BEGIN GLG6
\paragraph{GLG6. the good-orbit inequality from \(p_0\) to \(X'\).}
\label{claim:GLG6}
Under GLG4, let \(p_0\le p<X'\) be a good candidate prime.
The original quantities are
\[
 \begin{split}
 \Pr_v&=\sum_{s\in T_v\setminus V}(F_j^G(0,s)-b_s)+K_Pn'U,\\
 t_v&=S_{\rm kink}+K_1NN_1E_2
            +K_2N(Qc_6(n')+kQ+kl_p+c)+2n'c,\\
 D_p&=D_I(n').
 \end{split}                                                       \tag{GLG.20}
\]
Then the same \(\mathrm{StarAt}\) inequality (GLG.18) holds.

\emph{Proof.}
Use GLG3 to combine the first term of \(\Pr_v\) with
\(S_{\rm kink}\). Their sum is at most
\(\sum_{T_v\cap V}M_j+N(b'/x+\eta)\).
The term \(K_Pn'w\) is at most
\[
 NK_PG_0w\le NCw\,l/\delta.
 \]
The \(K_1NN_1E_2\) term is at most \(NClE_2\).
All remaining nonnegative terms, including the two ceiling
errors and the final \(1\), are at most
\[
 \begin{split}
 N\{&K_PG_0(\eta+c+N_m)
       +K_2(Qc_6(n')+kQ+kl_p+c)\\
       &+2G_0c+40C_{31}lE_I+3\}\le Nb_o.
 \end{split}
 \]
Here \(D_I(n')\le40n'\delta E_I+1\),
\(\lceil t_v\rceil\le t_v+1\), and (GLG.15) were used.
Thus
\[
 \Pr_v+\lceil t_v\rceil+D_p+1
 \le \sum_{T_v\cap V}M_j
       +N\{b'/x+\eta+Cwl/\delta+b_o+ClE_2\}.                      \tag{GLG.21}
\]
The original definition in this range is
\[
 \operatorname{lo}(C,b_o,p)
   =C(c+\eta)\log_p(2a)+b'/x+\eta+Cwl/\delta+b_o+ClE_2 .
 \]
Combine (GLG.16) and (GLG.21) and again use \(C\ge3\).
Every term, including the coefficient \(1\) of \(b_o\), is
unchanged. In particular no upper restriction on the \(b_s\)
was needed in the cancellation.
\hfill\(\square\)
% END GLG6

% BEGIN GLG7
\paragraph{GLG7. a single threshold proves the original \(\mathrm{RestStar}\).}
\label{claim:GLG7}
Define
\[
 \begin{split}
 c_G(\epsilon)&=\{(5+8\epsilon)10^{6r+1}+6\}(6r+1)2,\\
 c_N(\epsilon)&=18c_G(\epsilon)+2(5+8\epsilon)+2,\\
 C_*&=3+K_Tr+(13+4K_T+160C_{31})
                    +K_PC_{31}+K_1|c_N(\epsilon)|.
 \end{split}                                                       \tag{GLG.22}
\]
For every original family \(\mathrm{Fm}\), nonnegative \(K\),
\(C_{31}>0\) and real \(C_e,C_m\), choose \(C_0=C_*\).
For every \(C\ge C_0\), and every pair of original families
of instances \(X,Y\) satisfying the full
\(\mathrm{AdmFam}(\mathrm{Fm},K,X,Y,C_{31},C_e,C_m)\),
there is \(x_0\) such that, for every \(x\ge x_0\) and every
member \(i\), (GLG.18) holds simultaneously for every
original \(a_s,b_s\) in (GLG.2), every candidate prime,
every good orbit and every lift, with
\[
          b_o=b^{\,x^{-\epsilon/2}}=\exp(x^{\epsilon/2}).           \tag{GLG.23}
\]
This proves the literal \(\mathrm{RestStar}\), including the
order \(C_0\), then \(C\), then \(X,Y\), then \(x_0\).

\emph{Proof.}
Only \(\mathrm{adm1}\), \(\mathrm{adm5}(\epsilon/48)\),
the original parameter conditions and \(\mathrm{SizesAt}\)
are needed from the admitted record. The other fields are
retained in the quantified antecedent.
The definitions with their natural ceilings give, for \(x\ge e\),
\[
 G_p\le c_G(\epsilon)\log x\,x^{\epsilon/8},
             \qquad N_1\le c_N(\epsilon)\log x.                  \tag{GLG.24}
\]
Indeed the small-multiplicity ceiling is at most
\((5+8\epsilon)10^{6r+1}\log x+1\), and
\(\lceil x^{\epsilon/8}\rceil\le2x^{\epsilon/8}\).
Their product gives the first estimate. For the second,
the two arguments of the orbit-window maximum are
\(6eG_p\delta\) and \((5+8\epsilon)\log_2x\);
use \(e<3\), \(\log2>1/2\), and allow \(2\) for its
ceiling and added \(1\). This gives exactly \(c_N\).
Also \(k\le x^6\) gives
\(L_1(p)\le\lceil186\log x\rceil\) for every prime
(\(4(6+26)/\log2<186\)); hence \(L_2<514\log x\).

For completeness the only limiting facts used are, for
\(\alpha>0\) and fixed \(A\),
\[
 A\log x=o(x^\alpha),\qquad
 Ax^\beta=o(x^\alpha)\quad(\beta<\alpha).
 \]
Writing \(x=e^z\), the first follows, for example, from
\(e^{\alpha z}\ge(\alpha z)^2/2\); the second is
\(Ax^{\beta-\alpha}\to0\). Therefore one threshold gives
\[
 \begin{gathered}
 514\log x\le x^{\epsilon/8},\qquad
 3c_G\log x\le x^{\epsilon/8},\\
 4C_{31}\log x\le x^{1-\epsilon/8},\qquad x\ge e .
 \end{gathered}
 \]
The original sizes give \(\eta\le w\) and \(N_m,c\le w\).
Thus \(U\le4b'/x\), and these inequalities prove
\(G_0U\le b'\), \(c_0\le3b'\) for \(x\ge2\), and the
first five conditions in (GLG.13). Equations (GLG.22)--(GLG.24)
give all the stated lower requirements on \(C\).

Here is the full uniform estimate for the remaining terms.
Put \(\beta=23\epsilon/48\) and \(Z=\exp(x^\beta)\).
After increasing one threshold, \(Z\ge130\); all fixed
nonnegative constants used below are at most \(Z\); and
\[
 \begin{gathered}
 \eta,c,N_m,n,k,x^{13},\log x,\log_p a,G_0,N_1\le Z,\\
 40n'\delta\le Z\quad(n'\le n),\qquad Q_p\le Z\quad(p<X'),\\
 Q_p\le Z\quad(p<p_0).
 \end{gathered}                                                     \tag{GLG.25}
\]
Every assertion is uniform in \(i,p,n'\). Here are explicit
checks of the possibly non-polynomial ones. The original
sizes say \(n\le x^3\), \(k\le x^6\),
\(\log\eta\le x^{\epsilon/4}\), and
\(\log\Delta\le x^{\epsilon/4}/2\).
The inequality \(14\log x\le x^\beta\) bounds all powers
up to \(x^{14}\) by \(Z\). Since \(16k\le x^7\) for
\(x\ge40\), it gives \(c=16k\Delta\le Z\).
Next \(\log_pa\le2\log a\le x^2\),
\(G_0=C_{31}\log x\,x^{\epsilon/8}\le x^2\),
\(N_1\le x\), and \(40n'\delta\le x^4\) eventually.
Finally \(\mathrm{adm5}(\epsilon/48)\) gives
\[
 \log Q_p\le x^{11\epsilon/24+\epsilon/48}=x^\beta
       \quad(p<X'),
 \]
and its small-prime clause gives the still smaller exponent
\(\epsilon/8+\epsilon/48\) when \(p<p_0\).

It follows that \(kQ_p\le Z^2\) and
\[
         W_D\le Z^5+Z^7+Z^4\le Z^8 .
 \]
Choose the same threshold so that
\[
 14x^\beta\le x^\epsilon,\quad
 4x^{\epsilon/3}\le x^\epsilon,\quad
 13x^\beta\le x^{\epsilon/2}.
 \]
These are possible because \(\beta<\epsilon/2<\epsilon\).
They imply
\(Z^8\le x\exp(x^\epsilon-x^{\epsilon/3})=b'\) and
\(Z^{13}\le b_o\).
To check the entire expression (GLG.14), first use
\[
 D_I(n')\le
 Z\cdot Z\cdot Z^2(Z^5+Z^3)Z+1 .
 \]
Substitute (GLG.25) and this ceiling bound into \(B(n',p)\).
Its resulting polynomial is
\[
 8Z^{12}+40Z^{11}+8Z^{10}+40Z^9+Z^6
                     +6Z^4+5Z^3+19Z^2+3
          \le130Z^{12}\le Z^{13}\le b_o.                           \tag{GLG.26}
\]
This proves every condition of GLG4. Intersect the finitely many
scalar thresholds with the admitted-family threshold at
\(\epsilon/48\), and with the family's parameter threshold.
GLG5 and GLG6 then apply to all members and all window functions
at once. In particular the threshold has not been selected
after the orbit, lift, prime, or top weights.
\hfill\(\square\)
% END GLG7

% BEGIN GLG8
\paragraph{GLG8. from good-orbit inequalities to the original gain.}
\label{claim:GLG8}
Assume \(\mathrm{StarAll}(X,K,C,b_o)\), the original parameter
conditions, \(b_o\ge0\), \(C\ge3\), and
\[
 C\ge2K.\mathrm{dl1}+C_e',\qquad C\ge4K.\mathrm{dl2},\qquad
 e_p\le C_e'f k^2\log_pa\quad(p\le X'').                           \tag{GLG.27}
\]
The \(e_p=X.\mathrm{ep}(p)\) are the original nonnegative
integer exponents. Then for all \(a_s,b_s\) in (GLG.2), all
inert or signal primes \(p\le X''\), and all lifts,
\[
 0\le\mathrm{gain}_j(p)\le4|TT|a/p,                               \tag{GLG.28}
\]
and
\[
 \sum_{j=1}^m\sum_{s\in TTP(p)}
                 \{A_j(s)-\operatorname{lo}(C,b_o,p)\}
       -m\,\operatorname{fl}(C,p)
        \le\sum_{j=1}^m\mathrm{gain}_j(p)-\tau_p .                 \tag{GLG.29}
\]
Here every term is the original term:
\[
 \begin{split}
 \operatorname{fl}(C,p)&=CfN_m^2\log_pa+Cf(\log_pa)^2,\\
 \mathcal V_p&=\begin{cases}
      \{\text{good representatives}\},&
      p\text{ candidate and }\mathrm{Sig}(p)\le x^{-\epsilon},\\
      \varnothing,&\text{otherwise},
      \end{cases}\\
 \tau_p&=\begin{cases}
      0,&\mathcal V_p=\varnothing,\\
      e_p+m\sum_{v\in\mathcal V_p}(\lceil t_v\rceil+D_p+1),
                                  &\mathcal V_p\ne\varnothing,
      \end{cases}\\
 \mathrm{gain}_j(p)&=
      \sum_{v\in\mathcal V_p}\max(g_v^U-\Pr_v,(g_v^0)_+)
      +\sum_{v\in R_p\setminus\mathcal V_p}(g_v^0)_+ .
 \end{split}                                                       \tag{GLG.30}
\]
The last expression is replaced by zero if \(p\) is inert
and \(h_p=0\), or if \(p>X''\). The original \(TTP(p)\)
is all \(TT\) at signal primes or \(p\ge X'\), and otherwise
the tops whose representative lies in \(\mathcal V_p\).

\emph{Proof.}
The representatives partition \(HH\) into equal-length orbits.
Hence \(|R_p|n_p\le f\) (the signal convention multiplies
the orbit length by \(2\)), while \(\sum_{v\in R_p}N_v=|TT|\le f\).
Since \(\log_p(2a)\le2\log_pa\), summing (GLG.4) gives
\[
 \sum_{v\in R_p}d_v
 \le2K.\mathrm{dl1}\,fN_m^2\log_pa
                      +4K.\mathrm{dl2}\,f(\log_pa)^2 .
 \]
As \(k\le N_m\), (GLG.27) implies
\[
                     \sum_{v\in R_p}d_v+e_p
                                      \le\operatorname{fl}(C,p). \tag{GLG.31}
\]
Every \(\Pr_v\) is nonnegative. In the large-prime branch
its top terms equal \(o_j-\kappa_j\ge b'/x-S_c\ge0\);
in the small branch this is immediate. The other losses
are also nonnegative. Thus both \(g_v^0\) and \(g_v^U\)
are at most \(\gamma_p\sum_{T_v}\alpha_s\), and so are
the positive maxima appearing in the gain. Since
\(\gamma_p\le4/(3p)\) and \(\alpha_s\le6a/5\), summing
over the partition even gives the stronger upper bound
\(8|TT|a/(5p)\). This proves (GLG.28), including the
branches where the gain is defined to be zero.

If \(\mathcal V_p\ne\varnothing\), \(p\) is a candidate
and \(\mathcal V_p\) is the whole good-representative set.
Sum \(\mathrm{StarAt}\) over these representatives, and
then over the \(m\) lifts. The left side is the sum over
the actual \(TTP(p)\). The \(g_v^U-\Pr_v\) terms are
bounded by the selected gain maxima. The accumulated
rounding terms are exactly \(\tau_p-e_p\). Finally
\[
 e_p+m\sum_{\mathcal V_p}d_v
   \le m\bigl(e_p+\sum_{R_p}d_v\bigr)
   \le m\,\operatorname{fl}(C,p),
 \]
because \(m\ge1\) and \(e_p,d_v\ge0\). This proves (GLG.29).

Suppose next \(\mathcal V_p=\varnothing\). If \(TTP(p)\)
is empty, (GLG.29) follows from nonnegative gain and
\(\operatorname{fl}\). Otherwise \(p\) is signal or
\(p\ge X'\). At a signal prime the excluded terms in
\(\mathrm{g0}\) do not occur: an odd partner has first
coordinate \(1\), whereas an excluded partner is a top
with first coordinate \(0\). For \(p\ge X'\), every
excluded partner is a top and, for \(e\ge1\),
\[
 \frac{a_s}{p^e}
 \le\frac{a+a/(kx^{13})}{X'}
       =b'+\frac{b'}{kx^{13}}
       \le b'+b'/x\le F_j(\operatorname{partner}(p,s,e)).
 \]
Thus all excluded positive parts vanish. Apply GLG2 to
the remaining terms to obtain
\[
 \sum_{s\in T_v}\{A_j(s)-3c\log_p(2a)\}
                                      \le g_v^0+d_v .
 \]
The original \(\operatorname{lo}\) is at least
\(3c\log_p(2a)\): its base term has \(C\ge3\), and all
its other branches are nonnegative. Sum over all
representatives and lifts and use (GLG.31).
If in addition \(p\) is inert and \(h_p=0\), every partner
is \((0,s)\), the preceding top bound makes all \(A_j(s)\)
zero, and the inequality follows directly despite the
gain's special zero definition. These cases exhaust the
original prime quantifier.
\hfill\(\square\)
% END GLG8

% BEGIN GLG9
\paragraph{GLG9. the full original remaining counting theorem.}
\label{claim:GLG9}
For every original \(\mathrm{Fm}\), every nonnegative \(K\),
and every \(C_{31}>0\), \(C_e,C_m\ge0\), define
\[
 C_1=\max\{C_*,3,\,2K.\mathrm{dl1}+C_e+1,\,
                                      4K.\mathrm{dl2}\},
                  \qquad C_{\rm count}=C_1+1.                    \tag{GLG.32}
\]
This constant is chosen before the instance families \(X,Y\).
For every such pair satisfying the full original
\(\mathrm{AdmFam}\), there is a common \(x_0\) such that
for every \(x\ge x_0\), every member \(i\), and every pair
of integer top functions in (GLG.2), the exact original
\(\mathrm{GainBoundAll}(X_{x,i},K,C_{\rm count},
                 b^{x^{-\epsilon/2}})\) holds.
Thus the complete source endpoint
\(\mathrm{RemainingCountingStatement}(\mathrm{Fm})\)
has been proved.

\emph{Proof.}
GLG7 supplies \(\mathrm{StarAll}\) at \(C_1\).
The admitted record at the fixed auxiliary exponent \(1\)
supplies \(\mathrm{adm9}\) and \(N_m+1\le w\).
Its \(\mathrm{adm9}\) clause says
\(e_p\le(C_e+1)fk^2\log_pa\) for \(p\le a\).
Since \(w\ge500\), \(X''=a/w\le a\), so GLG8 applies.

It remains to replace \(\mathrm{askP}\) by the actual
\(\mathrm{askG}\), not to assume they are equal. By (GLG.7),
GLG1, and \(a_s\le2a\),
\[
           A_j(s)\le A_j^G(s)
                \le A_j(s)+N_m\log_p(2a).                         \tag{GLG.33}
\]
This also proves the full \(\mathrm{RequiresAll}\) assertion
for arbitrary integer \(a_s\le2a\), \(b_s\ge0\): when
\(a_s<0\) both positive-part sums are zero.
Increasing \(C_1\) by \(1\) increases
\(\operatorname{lo}\) by at least
\((c+\eta)\log_p(2a)\ge N_m\log_p(2a)\), and increases
\(\operatorname{fl}\). Each additional branch coefficient
is nonnegative; the \(b_o\), \(b'/x\) and \(\eta\) terms
whose coefficient is \(1\) remain unchanged.
Since \(TTP(p)\subseteq TT\), (GLG.33) converts (GLG.29)
term by term into exactly the \(\mathrm{askG}\) inequality.
The gain bounds are unchanged. Take the maximum of the
two family thresholds. This preserves every quantifier
in the stated endpoint and the original error (GLG.23).
\hfill\(\square\)
% END GLG9

\subsection{Actual finite places and the comparison lattice}

% BEGIN GLG10
\paragraph{GLG10. served indices, coverage, and the original \(S_v\) bound.}
\label{claim:GLG10}
On the actual degree functions \(a_L,b_L\), for every prime
and every served representative,
\[
                         S_v\le 2mN_va/p .                        \tag{GLG.34}
\]
The original served-set coverage condition holds identically.
For the original family, \(\mathrm{ServedGIdx}\) holds for every
member once
\[
              x\ge\max\{\mathrm{Fm}.x_1+1,\,
                                      4^{\,5/\mathrm{Fm}.ez}\}.   \tag{GLG.35}
\]

\emph{Proof.}
A served representative occurs only at a candidate prime, so
\(p<X'\). By definition
\[
 S_v=\left\lceil\sum_{j=1}^m
       \max\{0,g_v^U-\Pr_v-\max(g_v^0,0)\}\right\rceil .
 \]
GLG1 and the nonnegativity argument of GLG8 give
\[
 \max\{0,g_v^U-\Pr_v-\max(g_v^0,0)\}
       \le\gamma_p\sum_{T_v}\alpha_s\le\frac85N_va/p.
 \]
If \(N_v=0\), \(g_v^U=-d_v\le0\), and \(S_v=0\).
Otherwise \(N_v\ge1\), and \(a/p>b'\ge w\ge500\).
The ceiling costs at most \(1\), which is bounded by
\((2-8/5)mN_va/p\). This proves (GLG.34) with the
original coefficient \(2\), including an empty-top orbit.
If a prime is a candidate and its total deficit is at most
\(x^{-\epsilon}\), the definition in (GLG.30) makes its
served set the entire good-representative set. This is the
exact coverage assertion.

To check the local index, a nonempty served set implies a
candidate inert prime and hence a positive even Frobenius
order. The original arithmetic Frobenius identification
therefore gives the nonzero host component required by
\(\mathrm{GIdx}\). The original level-prime condition says
that level primes split, so the inert prime is not a level
prime. Under (GLG.35) the supplied split bound is at least
\(\lfloor x^{ez/5}\rfloor\ge4\). An inert prime below \(5\)
would be \(2\), which is in the specified small-prime set,
or \(3\), which splits by this bound; both are impossible.
Thus \(p\ge5\). Finally, if \(p\ne T_u\), primality of
\(p,T_u\) implies \(p\nmid T_u\). These are exactly the
fields of \(\mathrm{GIdx}\), for every word \(u\); the
threshold precedes all words and finite places.
\hfill\(\square\)
% END GLG10

% BEGIN GLG11
\paragraph{GLG11. the finite class sum and the original global loss.}
\label{claim:GLG11}
Assume the original \(\mathrm{ParametersFill}\) identifications,
including \(\Lambda=L.\mathrm{LaG}\). Then the actual finite
von Mangoldt class sum is exactly \(A_j(s)\), not merely
asymptotic to it. For every real \(C,b_o\), every prime \(p\),
and \(C'=4|C|+1\),
\[
 \operatorname{lo}(C,b_o,p)\le\operatorname{loG}(C',b_o,p),
 \qquad
 \operatorname{fl}(C,p)\le\operatorname{flG}(C',p).                \tag{GLG.36}
\]
Here the global loss is literally
\[
 \begin{split}
 \operatorname{loG}(C',b_o,p)
 ={}&C'(S_c+k+\eta)\log_p(2a)+P_G(p)+O_G(p),\\
 P_G(p)={}&
 \begin{cases}
 C'rb'\log x,&p\text{ inert},\ p<p_0,\\
 C'\{b'/x+\eta+S_c+(w+t)\log x/\delta\}+b_o,
                          &p\text{ inert},\ p_0\le p<X',\\
 0,&\text{otherwise},
 \end{cases}\\
 O_G(p)={}&
 \begin{cases}
 C'\log x\{\Lambda a/(pk)+ap^{-L_*(p)}\},
                          &p\text{ inert},\ p<p_0,\ p<X',\\
 C'\log x\{\Lambda a/(pk)+a/p^2\},
                          &p\text{ inert},\ p_0\le p<X',\\
 0,&\text{otherwise},
 \end{cases}\\
 L_*(p)={}&\left\lceil\log_p(k^2x^{78})\right\rceil+1,\qquad
 \operatorname{flG}(C',p)=C'fN_m^2\log_pa+C'f(\log_pa)^2 .
 \end{split}                                                       \tag{GLG.37}
\]

\emph{Proof.}
The supplied filling identity identifies the summand of
the finite class sum with
\((a_s/p^e-F_j(\operatorname{partner}(p,s,e)))_+\) for
odd \(e\), and with zero for even \(e\). Its range ends
at \(\operatorname{Nat.log}_p(a_s)+2\). Beyond this range
\(p^e>a_s\); since the filling is at least \(1\), every
omitted positive part is zero. The top degree is a positive
integer by (GLG.2), so the natural-log convention causes
no exceptional case.

For (GLG.36), \(C'\ge1\), \(C'\ge C\), and
\[
              C(4S_c+\eta)\le C'(S_c+k+\eta).
 \]
All multiplying logarithms are nonnegative.
The small-prime \(P_G\) term follows from \(C'\ge C\).
In the middle range, \(C'\ge1\) absorbs \(b'/x+\eta\),
\(C'\ge C\) absorbs \(Cw\log x/\delta\), and the
additional \(S_c\) and \(t\) terms are nonnegative.
The coefficient of \(b_o\) stays \(1\), even if \(b_o<0\).
When \(p<p_0\),
\[
 \log_p(k^4x^{104})-\log_p(k^2x^{78})
                       =\log_p(k^2x^{26})>1,
 \]
so \(L_*(p)\le L_1(p)\) and
\(ap^{-L_1(p)}\le ap^{-L_*(p)}\).
This proves the tail comparison; if \(C<0\), first use
\(C\le C'\) on its nonnegative tail, and then the
nonnegative coefficient \(C'\) on this exponent comparison.
The remaining prime branches are zero on both sides.
The formula for \(\operatorname{flG}\) proves its inequality
directly. This verifies the exact envelope for arbitrary
real \(C\), rather than imposing a new positivity hypothesis.
\hfill\(\square\)
% END GLG11

% BEGIN GLG12
\paragraph{GLG12. the actual local depth inclusion.}
\label{claim:GLG12}
Let \(w\mid p\) be served and \(p\ne T_u\), and use the
original \(\mathrm{LatCmp}\) and \(\mathrm{NuServedDef}\).
Then
\[
 \mathrm{LocLe}\!\left(w,\lfloor e_w\tau_p\rfloor,
                              \mathrm{NuLat}_u,M_{{\rm ram},u}\right)
                                                                    \tag{GLG.38}
\]
holds, where \(e_w=v_w(p)\). Consequently the exact
\(\mathrm{NuServed}\) predicate is supplied by these original
comparison data, with no rational-operator assumption.

\emph{Proof.}
Write \(K_0\) for the global coefficient field,
\(A=\mathcal O_{K_0,w}\), \(O\subset F\) for the comparison
DVR and its fraction field, \(\iota:K_0\hookrightarrow F\)
for the supplied embedding, and \(\phi\) for the supplied
semilinear comparison. By \(\mathrm{LatCmp}\),
\[
 \iota(z)\in O\ \Longleftrightarrow\ z\in A,\qquad
 M=O\phi(L_u),
 \]
and the images of a \(K_0\)-basis remain independent over \(F\).
These conditions imply the following descent fact:
if \(v\in V_u\) and \(\phi(v)\in O\phi(L_u)\), then
\(v\in A L_u\). To see this, choose an \(A\)-basis of
the finite free localization \(A L_u\). Its images are
independent over \(F\). The unique coordinates of
\(\phi(v)\) in this basis belong to \(O\), hence the
coordinates of \(v\) belong to \(A\). Clearing their
finitely many denominators produces \(s\in\mathcal O_{K_0}\),
\(s\notin w\), with \(sv\in L_u\).

Now take \(z\in\mathrm{NuLat}_u\).
The served defining identity gives
\(\phi(\nu_{\rm in}z)\in P\), where
\[
                  P=J\cap p^T M\cap U_F,\qquad
                  T=\max(0,\lceil\tau_p\rceil).
 \]
Thus \(\phi(\nu_{\rm in}(p^{-T}z))\in M\).
By the descent fact there is \(s\notin w\) such that
\(s\nu_{\rm in}(p^{-T}z)\in L_u\). It is already in the
rational ramified subspace, since \(\nu_{\rm in}\) has
that image. Hence \(sp^{-T}z\in M_{{\rm ram},u}\).
As \(p\in w^{e_w}\) and
\(\lfloor e_w\tau_p\rfloor\le e_wT\),
\[
          sz=p^T(sp^{-T}z)
                     \in w^{\lfloor e_w\tau_p\rfloor}
                                          M_{{\rm ram},u}.
 \]
This is the exact denominator-cleared definition of
\(\mathrm{LocLe}\). In particular the factor is \(e_w\),
not the valuation of \(p\) in the larger comparison DVR.
\hfill\(\square\)
% END GLG12

% BEGIN GLG13
\paragraph{GLG13. the original local operators preserve both lattices.}
\label{claim:GLG13}
For every original \(G=GAtP(p)\) and each local vertex \(v\),
\[
                        y_vM_{\rm na}\subseteq M_{\rm na},
                  \qquad y_vM\subseteq M .                        \tag{GLG.39}
\]
These are the two exact fields of
\(\mathrm{FormalGroupSystem.adm12}\), for the original
dual-Weyl and Weyl integral coefficient models.

\emph{Proof.}
We first verify the elementary integral step, including the
Weyl model for which preservation cannot be inferred merely
from containment in the dual-Weyl model. Let \(\lambda\) be
decreasing of length \(q\), \(e=\lambda_q\), and
\(\nu_i=\lambda_i-e\ge0\). Let \(z\ne0\) in \(O\), and
\[
                 D_r(z)=\operatorname{diag}
                         (\underbrace{z^{-1},\ldots,z^{-1}}_r,
                                           1,\ldots,1).
 \]
The tableau basis in CB2 is homogeneous for row degrees.
In every one of its monomials the total degree in the
first \(r\) rows is at most \(N=\sum_{i\le r}\nu_i\):
each tableau column has distinct row entries, so it uses
at most \(r\) of those entries, and summing over columns
gives \(\sum_i\min(r,\text{column height}_i)=N\).
Thus the corresponding diagonal substitution, multiplied
by \(z^N\), preserves the integral dual-Weyl module.

For the Weyl module use its defining integral generators:
words in row divided-power operators applied to the
highest-weight product of leading minors. Each generator
is homogeneous in row degrees. Indeed the coefficient
of \(T^h\) in a row substitution changes its two row
degrees by \(+h,-h\). The generators lie in the dual-Weyl
module: generic column invariance and column multidegree
are preserved by row substitutions, as they commute
with the column substitutions; comparing coefficients
of \(T\) proves this for divided powers. Consequently a
nonzero homogeneous generator has first-\(r\)-row degree
\(d\le N\), by the preceding support bound. Multiplication
by \(z^N\) followed by \(D_r(z)\) is exactly multiplication
of this generator by \(z^{N-d}\in O\). It therefore
preserves the Weyl module itself.

In the original signed-weight lattice the determinant
factor is \(\det(g)^{-e}\). Replacing \(g\) by \(gD_r(z)\)
adds the factor \(z^{re}\). Since
\[
                 \sum_{i\le r}\lambda_i=N+re,
 \]
the resulting statement for either model is precisely
\[
 z^{\sum_{i\le r}\lambda_i}L_\lambda(\operatorname{col}_O g)
       \subseteq L_\lambda(\operatorname{col}_O(gD_r(z))).        \tag{GLG.40}
\]
There is no assumption that \(e\) is nonnegative.
If two basis matrices give the same lattice, they differ
by a matrix in \(\mathrm{GL}_q(O)\). The dual-Weyl model
is stable under row changes. The Weyl model is stable too:
elementary unipotent row changes are finite sums of the
defining divided powers, diagonal units act by integral
units on homogeneous generators, and these generate
\(\mathrm{GL}_q(O)\) by elimination over a local ring.
The determinant correction is a unit even for negative
\(e\). Thus (GLG.40) holds for the actual basis-independent
models \(H_0\) and \(\mathrm{Dl0}\), as well as for a
chosen basis.

For an original edge \(\Lambda\to\Lambda'\), the elementary
divisor theorem over its unramified coefficient DVR gives
a basis in which the target is obtained by multiplying
the first \(r\) basis vectors by \(p^{-1}\). Indeed
\(p\Lambda'\subseteq\Lambda\subseteq\Lambda'\) and the
specified quotient dimension is \(r\); the elementary
divisors are therefore \(p\) in exactly \(r\) positions
and \(1\) elsewhere. In the actual comparison ring the
specified element \(z\) satisfies \(z^2=p\). Each Frobenius
twist fixes \(p\), so the extended edge consists of two
successive steps \(D_r(z)\), with the original midpoint
between them. Apply (GLG.40) twice in every tensor factor.
The product of the factors is
\[
 \prod_s z^{2\sum_{i\le r}\lambda_{s,i}}
                         =p^{\sum_s\sum_{i\le r}\lambda_{s,i}}
                         =p^{G.P.\mathrm{ArV}(v)}.               \tag{GLG.41}
\]
This is exactly the normalization of the original \(y_v\).

Finally \(y_vF\) at a lattice tuple is this factor times
the finite sum of \(F\) over incoming neighbors in the
\(v\)-coordinate. Equation (GLG.41) places each summand
in the target tensor lattice; all other coordinates stay
fixed. If a tuple is invalid, each possible incoming
term is zero, since an incoming edge has lattice endpoints
and cannot repair an invalid different coordinate.
The edge relation is equivariant under the given level
group, so this sum preserves its prescribed character
equivariance. These are all three defining conditions of
the actual function lattice, for both coefficient models.
This proves (GLG.39) over the given DVR, without extension
or change of the rank or of \(\mathrm{ArV}\).
\hfill\(\square\)
% END GLG13

% BEGIN GLG14
\paragraph{GLG14. coefficient rationality supplies the literal \(\mathrm{NPDesc}\).}
\label{claim:GLG14}
For the original \(\mathrm{LatCmp}\), retain the original
\(\mathrm{GatTies}\), and assume the original
\(\mathrm{ValRat}\): at each genuine lattice tuple every
coefficient of \(\phi(v)\) is in \(\iota(K_0)\).
Then every original \(y_v\) descends to a \(K_0\)-linear
operator on the same global vector space. Under the
original \(\mathrm{VramDef}\) this proves
\(\mathrm{NPDesc}\) for the original selected lattice.
In particular the original
\(\mathrm{GroupCompatibilityFamily}\) supplies eventual
\(\mathrm{NPDescF}\); its comparison record is unchanged.

\emph{Proof.}
Here is the required rational descent, with no Galois
assumption on \(F/K_0\). Choose a \(K_0\)-linear retraction
\(\pi:F\to K_0\) of \(\iota\), by extending \(1\) to a
\(K_0\)-basis of \(F\). Polynomial coefficient evaluations
at all component indices and lattice tuples separate
the ambient function space. Suppose
\(H=\sum_i a_i\phi(v_i)\) lies in the \(F\)-span of the
comparison image and all those evaluations are rational.
Set \(v=\sum_i\pi(a_i)v_i\). If an evaluation of
\(\phi(v_i)\) is \(\iota(r_i)\), its value on \(\phi(v)\)
is \(\iota(\sum_i\pi(a_i)r_i)\); this is
\(\iota\pi(\sum_i a_i\iota(r_i))\), equal to the
rational value on \(H\). Separation gives \(\phi(v)=H\).
The comparison is injective by its stipulated independent
basis images, so this preimage is unique.

The apparent restriction in \(\mathrm{ValRat}\) to genuine
tuples causes no gap. Every element of \(M\) vanishes
at invalid tuples by definition. So does its \(F\)-span.
Since \(M=O\phi(L_u)\) and \(L_u\) spans the global space,
\(\phi(v)\) belongs to that span and vanishes there as well.
Thus all coefficient evaluations used above are rational.
By GLG13, \(y_v\) preserves the \(F\)-span of \(M\).
Its neighbor-sum formula and its rational scalar
\(p^{\mathrm{ArV}(v)}\) preserve rational coefficients.
The preceding argument therefore gives a unique preimage
of \(y_v\phi(v')\) for every \(v'\). Uniqueness and
linearity of \(y_v,\phi\) make this a \(K_0\)-linear
operator. This proves the entire \(\mathrm{YvRat}\)
predicate, including every local vertex.

Let \(\mathcal V\) be the actual served set and set
\[
 s(B)=\sum_{v\in\mathcal V\setminus B}S_v,\qquad
 T=\max(0,\lceil\tau_p\rceil),\qquad
 J=\sum_{B\subseteq\mathcal V}
           p^{s(B)}\left(\prod_{v\in B}y_v^m\right)M .            \tag{GLG.42}
\]
The \(S_v\) are already nonnegative integers. Products
use the very order in the original finite-set list; no
commutation assertion is necessary for rational descent.
The descended operators show that \(J\) is the
\(O\)-span of the images of rational vectors
\(p^{s(B)}\prod_{v\in B}(y_v^{\rm rat})^m z\), \(z\in L_u\).
Likewise \(p^TM\) is the extension of \(p^TL_u\), and
\(\mathrm{VramDef}\) identifies \(U_F\) with the extension
of the rational ramified subspace.

The original \(\mathrm{GatTies}\) supplies the fraction-field
and local coefficient-ring identifications. Together with
\(\mathrm{LatCmp.embW\_int}\), it gives the injective local
map \(A\to O\) used here. This map is faithfully flat:
it is torsion free over the DVR \(A\), hence flat, and
its maximal ideal remains proper. Under the independent
basis comparison, tensoring therefore commutes with
finite sums, images, and intersections, the last by
the kernel sequence for \(B_1\oplus B_2\to B_1+B_2\).
It also commutes with intersection with a rational
subspace, using its quotient vector space and the same
kernel argument. It follows that
\[
                   P=J\cap p^TM\cap U_F
 \]
is the extension of its rational preimage. This is the
literal \(\mathrm{NPDesc}\), not an identification of
arbitrary commensurable lattices.

Finally \(\mathrm{GroupCompatibilityFamily}\) includes
the actual \(\mathrm{CmpTied}\) polynomial-value record:
at a lattice tuple its \(\mathrm{val\_cmp}\) identity
writes \(\phi(v)\) as coefficient extension of the
specified polynomial over \(K_0\). Thus it supplies
exactly \(\mathrm{ValRat}\). Its \(\mathrm{LatDefs.vram}\)
is the \(\mathrm{VramDef}\) just used. Taking its common
eventual threshold proves \(\mathrm{NPDescF}\) for the
same family. Existence of this compatibility record is
an explicit antecedent, not a consequence of rational
descent.
\hfill\(\square\)
% END GLG14

% BEGIN GLG15
\paragraph{GLG15. finiteness and the exact normalized length.}
\label{claim:GLG15}
Retain the actual comparison, \(\mathrm{GatTies}\),
\(\mathrm{LatDefs}\), \(\mathrm{NuDef}\), and
\(\mathrm{NuUnserved}\), and use GLG14 for \(\mathrm{NPDesc}\).
Then the original \(\mathrm{LatticeLengthsAtPlace}\)
holds at every finite place. In particular at every
served \(w\) with \(p\ne T_u\) there is an integer \(l\ge0\)
such that
\[
 \operatorname{length}_O((M\cap U_F)/P)=l,\qquad
              l=(L.\mathrm{lenMN}(u,w))\,v_O(p).                 \tag{GLG.43}
\]
At an unserved place with \(p\ne T_u\),
\(L.\mathrm{lenMN}(u,w)=0\).

\emph{Proof.}
The original full global column lattice \(L_u\) is finite;
its intersection with the rational ramified space is
finite over the noetherian integer ring. The comparison
and flat intersection calculation identify \(M\cap U_F\)
with its extension, so it is a finite \(O\)-module.
Let \(s=\sum_{v\in\mathcal V}S_v\). The \(B=\varnothing\)
term in (GLG.42) gives \(p^sM\subseteq J\); hence
\[
                  p^{s+T}(M\cap U_F)\subseteq P .                \tag{GLG.44}
\]
The element \(p\) is nonzero and is a nonunit of \(O\).
Indeed its image comes from the rational prime of the
original unramified DVR through an injective local map,
as required by \(\mathrm{GatTies}\). A finite module
annihilated by \(p^{s+T}\) has finite length over a DVR.
This proves finiteness rather than assuming the local
length to be an integer.

Now apply the proved comparison SR2 to these very lattices,
using GLG14 and the served defining identity.
If \(a_w=v_O(\pi_w)>0\), scalar extension multiplies each
Smith exponent by \(a_w\), while \(v_O(p)=a_we_w\).
The source's \(\mathrm{NuDef}\) normalizes
\(L.\mathrm{lenMN}(u,w)\) as the \(A\)-length divided
by \(e_w\). Consequently
\[
 \operatorname{length}_O((M\cap U_F)/P)
   =a_w\operatorname{length}_A((M_{{\rm ram},u})_w/
                                             (\mathrm{NuLat}_u)_w)
   =v_O(p)\,L.\mathrm{lenMN}(u,w),
 \]
which is (GLG.43). The basis and ramified-space
identifications are unchanged, so their ranks agree too.
If \(w\) belongs to the original support, has the original
\(\mathrm{GIdx}\) guard, and its served set is empty,
(GLG.42) has only the empty-subset term:
\(J=M\) and \(\tau_p=0\). Thus \(P=M\cap U_F\), and
the local quotient has length zero. The guard excludes
\(p=T_u\), so the original
\(\mathrm{NuUnserved}\) also gives
\(L.\mathrm{lenMN}(u,w)=0\). This proves the first
clause of \(\mathrm{LatticeLengthsAtPlace}\) in its
empty-served case as well.
At an unserved regular place the original
\(\mathrm{NuUnserved}\) gives equality of the two localized
global lattices, hence length zero. The place \(p=T_u\)
keeps its separate original definition and is not
asserted to satisfy this regular-place zero conclusion.
\hfill\(\square\)
% END GLG15

% BEGIN GLG16
\paragraph{GLG16. the actual selected tops give the full second admissibility field.}
\label{claim:GLG16}
Suppose the original family satisfies \(\mathrm{TopsF}\):
\(\mathrm{ChosenTT}\), and \(\mathrm{ChosenBad}(C_b)\)
for one \(C_b>0\). Then the literal \(\mathrm{Adm2F}\)
holds, on the selected original subset \(TT\). If
\(x_a,x_b\) are the two supplied thresholds, one may take
\[
 x_0=\max\{x_a,x_b,\mathrm{Fm}.x_1+1,
                                     (16C_b+2)^{3/\epsilon}\}.    \tag{GLG.45}
\]

\emph{Proof.}
The chosen-top cardinality is already
\(\delta n/2\le|TT|\le2\delta n\).
Put \(T=x^{\epsilon/3}\). For any set \(\mathcal P\)
of positive integers less than \(\lceil e^T\rceil\),
\[
 \sum_{p\in\mathcal P}\frac{\log p}{p}
 \le2T\sum_{j=1}^{\lceil e^T\rceil}\frac1j
 \le2T\{1+\log\lceil e^T\rceil\}\le2T(T+2).                       \tag{GLG.46}
\]
The first step uses \(\log p\le T\) and
\(1/p\le2/(p+1)\); the harmonic bound follows by
integrating \(1/t\), and
\(\lceil e^T\rceil\le2e^T\) gives the last step.
In particular this applies to the original finite set
of large-order inert primes below \(X'\).
The \(\mathrm{ChosenBad}\) inequality says
\[
 \sum_{p\in\mathcal P}\frac{\log p}{p}
                       |\mathrm{badTops}(p)|
    \le4C_bx^{-2}\delta n
                    \sum_{p\in\mathcal P}\frac{\log p}{p}.
 \]
The actual \(\mathrm{BTT}\) sum is over a subset of
this set and has nonnegative summands. Therefore it is
at most \(8C_bx^{-2}\delta nT(T+2)\).
The threshold makes \(T\ge16C_b+2\), and hence
\(16C_b(T+2)\le T^2\). Thus
\[
 \mathrm{BTT}\le(\delta n/2)x^{-2}T^3
                      \le |TT|x^{-2+\epsilon}.
 \]
This is precisely the other clause of \(\mathrm{adm2}\).
The proof uses no asymptotic prime-counting theorem and
does not resample the selected top set.
\hfill\(\square\)
% END GLG16

\subsection{The original global installation and its remaining inputs}

% BEGIN GLG17
\paragraph{GLG17. the full literal finite-place installer.}
\label{claim:GLG17}
Fix the original \(D,L,X,Y,K\), with \(K\ge0\), and retain:
the actual \(\mathrm{GatTies}\), \(\mathrm{LatTied}\),
\(\mathrm{ParametersFill}\), \(\mathrm{NuDef}\),
the support condition containing every \(w\) with \(p\le X''\),
the actual length identifications of GLG15, and the following
two local-growth inputs:
\[
 \begin{gathered}
 \mathrm{MeanKappaBound}(X,K),\\
 \mathrm{admG1}(D,L,X,Y,K).
 \end{gathered}                                                     \tag{GLG.47}
\]
The latter is the original local-lattice assertion; its
length clause, on the same \(M,P\), is
\[
 \operatorname{length}_O((M\cap U_F)/P)
       \le v_O(p)\,d_u(\tau_p+\mathrm{io}_p).                     \tag{GLG.48}
\]
All its original additional clauses remain in the antecedent.
Assume \(\mathrm{ServedGIdx}\), and the original counting
conclusion at \(C,b_o\), with the ask comparison of GLG9.
Then
\[
                    \mathrm{LinkG\_at}(X,K,4|C|+1,b_o)           \tag{GLG.49}
\]
holds with all five original conjuncts.

\emph{Proof.}
At a served place with \(p\ne T_u\), GLG12 gives the
depth inclusion. GLG15 identifies the integer in
(GLG.48) with \(v_O(p)L.\mathrm{lenMN}(u,w)\).
The local-model rank is exactly \(d_u=D.\mathrm{duS}(u)\)
by \(\mathrm{GatTies}\), and \(v_O(p)>0\) by GLG15.
Division proves
\[
           L.\mathrm{lenMN}(u,w)\le d_u(\tau(w)+\mathrm{io}(w)).
 \]
The filling record identifies \(\tau,\mathrm{io}\)
with the original \(\tau_p,\mathrm{io}_p\).
If the served set is empty these are both zero by their
definitions, and the length is zero by GLG15. The depth
inclusion then has exponent zero and is exactly the
original regular-place inclusion in \(\mathrm{NuDef}\).
This proves the entire first conjunct, including its
empty-served-set implication.

GLG10 proves the second conjunct with constant \(2m\),
and proves the fourth conjunct, the exact coverage.
The fifth conjunct is the unchanged first input in
(GLG.47); GW2 supplies this input whenever its actual
canonical-deficit hypotheses are realized.

For the third conjunct take any signal or inert finite
place with \(p\le X''\). The support condition activates
the filling record's gain identity. Its top-subtype
bijection identifies \(L.\mathrm{TTw}(w)\) with the
original \(TTP(p)\); its coefficient identities
identify \(L.\mathrm{gain}(j,w)\) and \(\tau(w)\)
with \(\mathrm{gain}_j(p)\) and \(\tau_p\).
GLG11 identifies the finite class sum with \(A_j(s)\).
The inequality \(A_j(s)\le A_j^G(s)\), followed by
(GLG.36), makes the required global left side no larger
than the left side of the original counting conclusion.
The nonnegative factor \(m\) preserves the direction
when subtracting \(\operatorname{flG}\).
The gain bounds are exactly (GLG.28), so this proves
the whole third conjunct. Every conjunct concerns the
same supplied member, lattices, filling, degree functions,
and words; their witnesses have not been chosen separately.
\hfill\(\square\)
% END GLG17

% BEGIN GLG18
\paragraph{GLG18. the uniform family consequence and the exact frontier.}
\label{claim:GLG18}
Let \(\mathrm{Fm}\) be the original family. Suppose its
actual \(K=\mathrm{Fm}.GK\), \(X=\mathrm{Fm}.Gx\),
\(Y=\mathrm{Fm}.Gy\) satisfy the full original
\(\mathrm{AdmFam}\) for fixed \(C_{31}>0,C_e,C_m\ge0\);
retain \(\mathrm{GroupCompatibilityFamily}\) (or the
more precise \(\mathrm{NPDescF}\) it supplies), and retain
the exact eventual \(\mathrm{MeanKappaBound}\) on these
same instances. The \(\mathrm{admG1}\) field of the
admitted record supplies the original
\(\mathrm{FamilyFirstGrowthBound}\); when admissibility
is assembled from separate producers this field must
be supplied as well.
Then the complete literal \(\mathrm{LinkG}(\mathrm{Fm})\)
holds. One can use
\[
 C_G=4C_{\rm count}+1,\qquad
                   \delta_G(x)=x^{-\epsilon/2}.                  \tag{GLG.50}
\]
This is a proved implication on the original family;
it does not assert the existence of its retained input records.

\emph{Proof.}
GLG9 supplies one \(C_{\rm count}\ge0\) and the counting
inequalities at \(b^{\delta_G(x)}\), uniformly in the
member and all later choices. In particular they hold
at the actual \(a_L,b_L\), which satisfy (GLG.2).
The mean-deficit hypothesis is exactly the endpoint
proved in GW2 under its recorded original hypotheses:
the canonical \(\kappa_I\) identity in \(\mathrm{Fm.ex}\),
GI8 with its fixed coefficient bounds, and the small-prime
coefficient estimate. Those hypotheses are retained
when GW2 is used to discharge this input. No mean bound
for a new or independently selected deficit is being used.

The original family record already supplies, beyond
its starting threshold, \(\mathrm{GatTies}\), the actual
comparison and \(\mathrm{LatDefs}\), the filling identities,
\(\mathrm{LatTied}\), \(\mathrm{NuDef}\), the unserved
identity and the support condition. GLG14 supplies
\(\mathrm{NPDescF}\) from the stated compatibility
antecedent, and GLG15 supplies the required actual lengths.
GLG10 and GLG12 supply the finite-place conditions.
Take the maximum of their finitely many thresholds,
the admitted-family thresholds used in GLG9, the
stated mean-bound threshold and the admitted
first-growth threshold.
The constant \(C_G\) and function \(\delta_G\) are
already fixed before every later \(x,i,u,w,j\).
Apply GLG17 to every member at this common threshold.
Since \(\epsilon>0\), \(\delta_G(x)\to0\), giving
all quantifiers of the original \(\mathrm{LinkG}\).
\hfill\(\square\)

The finite local projection certificates are not a new
frontier in this statement: TC6 and EA9, specialized by
GI8/GII8 and transported by BC8 and GJ4, already prove them
under their original actual-model hypotheses. GLG13
adds the literal integral-operator supplier, GLG14
its rational descent, GLG16 the exact selected-top
admissibility consumer, and GLG9 the full previously
unexpanded logarithmic counting theorem. What remains
to assert the whole global leaf is the realization,
on this same family, of the retained compatibility,
the full admitted inputs, and the local first-growth
assertion (or its already reconstructed simultaneous
exceptional-direction inputs). No conditional estimate
has been relabeled as a construction of those data.
% END GLG18
